\section{The Effectiveness Border of Distillation Algorithms}

When OPD algorithm was initially introduced, its superior performance was believed to come from its on policy nature~\citep{thinking_machines_lab_-policy_2025}. SFT trains the student on the trajectory generated by the teacher. Even if the student learns well on prediction over teacher's prefix, if it can not generate a prefix with the same quality, the performance will still be limited. OPD method however, directly uses teacher's predictive distribution to correct the student on its own trajectory, and avoids this drawback~\citep{thinking_machines_lab_-policy_2025}. However, this discussion ignored another important impact in OPD: The teacher's probability prediction on the student's prefix may not be sufficiently reliable as well. In some previous works we did observe some cases where OPD may face stability issue and failed to improve the student model~\citep{li_rethinking_2026}. Clearly this issue does not exist in SFT.

It is difficult to determine precisely which of the two issues discussed above has a dominant impact on performance. Beyond such relatively subjective analysis, however, we aim to establish rigorous theoretical results and characterize the regimes in which the two algorithms are effective. Since all distillation methods have the same target distribution, one primary theoretical difference may lie in the estimation of the gradient. This leads to some interesting results.

\subsection{Proximal Behavior}

We study the proximal regime, where the student distribution is close to the teacher. At the \emph{proximal limit}, they coincide, while the teacher is held fixed during differentiation. Both negative-KL population objectives and their gradients then vanish, but their sampled gradient estimators behave differently.

\Needspace{11\baselineskip}
\begin{lemma}[Gradient noise at the proximal limit]
\label{lem:proximal-gradient-noise}
Using $m$ independent responses to a prompt $x$. At
$\pi_{\theta_0}(\cdot\mid x)=\pi_t(\cdot\mid x)$, both gradient estimators are unbiased with mean zero, but
\begin{equation}
 \widehat g_{\mathrm{OPD}}(\theta_0)=0\quad\text{almost surely},
 \qquad
  \operatorname{Cov}\bigl(\widehat g_{\mathrm{SFT}}(\theta_0)\bigr)
 =\frac1m F_x(\theta_0),
 \label{eq:proximal-gradient-noise}
\end{equation}
where $F_x(\theta_0)$ is the student Fisher information at prompt $x$.
Thus SFT typically has nonzero gradient noise at, whereas OPD vanishes for every sampled batch.
\end{lemma}
\noindent\hyperref[app:proximal-gradient-noise]{Proof: Appendix~\ref*{app:proximal-gradient-noise}.}

Under the smoothness assumptions, this strict gap still persists in a neighborhood of $\theta_0$: OPD has lower mean-squared error in estimating its population gradient than SFT has. This establishes a local advantage in gradient estimation in a range of proximal case, and could provide a theoretical explanation on the superior performance of OPD.

\subsection{Distal Behavior}

The distal regime can exhibit different behavior. When the student assigns appreciable probability to responses that the teacher considers very unlikely, the log-ratio weight $w_\theta(y\mid x)=\log\pi_t(y\mid x)-\log\pi_\theta(y\mid x)$ can amplify OPD's gradient noise. Under specific conditions, OPD's gradient estimation variance can be proved to exceed that of sampled-response SFT by an arbitrarily large factor, see \hyperref[app:distal-gradient-variance]{Appendix~\ref*{app:distal-gradient-variance}}.

% \begin{takeaway}
% We theoretically show that OPD provides more accurate gradient estimates than SFT in the proximal regime, whereas the opposite holds in the distal regime. This finding partially explains their differences in performance and stability across these regimes.
% \end{takeaway}

% \subsection{Empirical Validation}

%The preceding analysis establishes a local gradient-estimation advantage for OPD and a conditional variance comparison in the distal regime.

% We conduct a series of experiments to examine how these distinctions relate to distillation behavior. Experimental settings, training curves, and evaluations are provided in 
These results are also supported by our experiments, see Appendix~\ref{app:distillation-experiments} for more details.
